\documentclass[11pt]{article}
\usepackage[margin=1in]{geometry}
\usepackage{amsmath,amssymb,amsthm,booktabs,hyperref}
\newtheorem{proposition}{Proposition}
\title{The supplied third Janko geometries without a distinguished action}
\author{A finite application of the accepted incidence certificates}\date{}
\begin{document}\maketitle
Let $G=J_3$, of order $N=50\,232\,960$, and let $\Gamma_i$, $1\leq i\leq21$,
be the literal rank-three coset geometries in the acquired archive.
Indices refer to its final assignment, which exchanges construction entries
$8$ and $9$. The accepted eligibility certificates prove that all twenty-one
are firm, residually connected, flag-transitive, residually weakly primitive,
locally doubly transitive and satisfy $\mathrm{IP}_2$.
An isomorphism below may permute the types and need not preserve the
specified image of $G$.

\begin{proposition}
The supplied geometries have nineteen isomorphism classes.
The only identifications are $\Gamma_{13}\cong\Gamma_{14}$ and
$\Gamma_{16}\cong\Gamma_{17}$.
\end{proposition}
\begin{proof}
Write $G_J=\bigcap_{j\in J}G_j$ for a base flag stabilizer, with
$G_\varnothing=G$. Flag transitivity implies that the number of flags
of type $J$ is $N/|G_J|$: the orbit map sends a right coset $G_Jg$ to the
translate of the base flag, and equality of translates is precisely equality
of these cosets. Therefore the multiset of all eight numbers $N/|G_J|$
is intrinsic to the incidence geometry. A type permutation merely permutes
the eight subsets. Since $N$ is the same for every input, the corresponding
multiset of eight stabilizer orders is also an intrinsic separating invariant.
This argument does not assume that an arbitrary isomorphism normalizes $G$.

The acquired intersection orders give sixteen distinct profiles. Their
only nonsingleton fibres are
\[
 \{7,8\},\qquad \{9,10,11\},\qquad
 \{13,14\},\qquad \{16,17\}.
\]
The saved whole-family automorphisms, conjugators and type permutations
give actual incidence isomorphisms for the last two pairs: a map carrying
the subgroup family to its target induces a well-defined bijection of
cosets and preserves nonempty intersections in both directions.

The accepted common point-graph argument separates $7$ from $8$ and
$9$ from $11$ even without a distinguished group action. It proves that
any such geometry isomorphism must normalize the point action of $G$:
the point graph has full automorphism group of order $2N$, whose derived
subgroup is the distinguished $G$, and the other objects have distinct
point shadows. The incompatible involution counts in the resulting
parabolics then exclude both isomorphisms. We use that complete argument,
not merely its invariant under group automorphisms.

Finally $\Gamma_{10}$ satisfies IP, whereas $\Gamma_9$ and $\Gamma_{11}$
fail WIP. Both properties are expressed solely through flags, incidences
and shadows, so are preserved by arbitrary isomorphisms and by type
permutations. The implication IP $\Rightarrow$ WIP separates $10$ from
both other members of its fibre. Thus the first repeated fibre splits
into two classes and the second into three, while each of the last two
is one class. The number is $16+1+2=19$.
\end{proof}

\section*{Source use and the scope of the count}
The seventeen diagram labels of the acquired Janko supplement have
twenty-one printed occurrences. Labels $3.7$ and $3.8$ each account for
two distinct classes; labels $3.11$ and $3.13$ each account for one class
despite their two supplied occurrences. The other thirteen labels each
account for one. In particular, counting the seventeen labels is not
counting the nineteen supplied incidence classes.

L091 is D.~Leemans, \emph{The residually weakly primitive geometries of HS},
Australasian Journal of Combinatorics \textbf{33} (2005), 231--236.
Its introduction expressly treats the Janko classification as previous
work, citing reference~[9]: D.~Leemans, \emph{The residually weakly primitive
geometries of $J_3$}, Experiment. Math.\ \textbf{13} (2004), 429--433
(our historical identifier L097). Its Table~1 in Section~3 quotes
comparison counts from those earlier classifications. The new HS
classification and the two Hoffman--Singleton constructions are separate
claims of L091, supported by their own accepted proofs.

The proposition proves the exact quotient of the supplied Janko roster
and, unconditionally on any global census, the existence of at least
nineteen eligible rank-three geometries. It does not enumerate all Janko
subgroup families. To conclude that nineteen is the global rank-three
total, one must retain the earlier classification's coverage premise,
in the explicit form that every eligible geometry occurs in this roster.
To interpret all nineteen as nontruncations, one must also retain its
absence-of-higher-rank premise. These are borrowed classification inputs,
not conclusions of the proposition; the old numerical value seventeen
is not imported as a premise. In particular no claim of an independently
reproved global Janko census is made.

The full text of L097 is unavailable in the fixed accessible catalogue,
and L097 is outside the active 105 mathematical originals. The acquired
thirteen-page author supplement supplies auxiliary data and does not
change that status. L096, whose supplied group families and intersection
certificates are reused here, is an active original with an accepted
corrected principal recovery. Neither use awards a second original or
recovery credit to L097. The Janko supplement's other diagram metrics,
abstract subgroup names and full automorphism groups are not new
assertions of L091 and need not be reclassified to audit its comparison
table. Existing positive certificates for some of them remain available.

The read-only verifier acquires the accepted finite inputs, reconstructs
every flag-count profile and this quotient, and records the historical
source correspondence separately. It performs no new group search.
The accepted GAP and Traces calculations retain their stated software
premises. This packet makes no Lean, conjecture or compression claim.
\end{document}
